\justifying
\NSFCBodyText

\subsubsection{项目特色}

项目以教育任务中的实际支持边界为研究起点，将教学适切性、证据忠实性与行动授权放在同一行为过程内考察。课程答疑和文献学习辅助共用可复现任务与工具环境，使风险分类、过程评价和干预后果能够相互核对。相同最终答案对应不同轨迹、相同行动对应不同教学要求的受控案例，有助于辨认最终结果难以揭示的风险，也避免把所有教育问题统一解释为模型能力不足。

研究规模服从一年期开放课题条件，以已有模型、公开材料和沙箱为主要载体。成果围绕同一套任务、日志和证据规则形成，评测基准、监测原型与案例报告能够互相追溯。对于真实场景条件尚未落实的部分保留授权前提，使核心机制检验不依赖高成本的全校部署。

\subsubsection{创新点}

\textbf{教育情境约束与过程证据共同作用的可信评测。}通用风险轨迹基准和近期教学工作流评测已经提供了重要基础。本项目拟进一步以允许帮助程度、任务目标和授权范围的受控变化，检验同一行为风险含义的改变，并在等信息通用评价对照下分析教育约束与过程证据的交互作用。创新的判断依据是跨任务家族的增益及其消融反证，而非仅增加教育题目或评测维度。

\textbf{以时间窗口和教学代价共同检验监测的有效性。}围绕首次可观察证据、首次告警和风险动作提交，区分事后定位与动作前发现，再在相同初始任务上追踪执行、重规划、阻断和人工复核产生的分支结果。研究拟说明哪些风险值得自动干预、何时需要保留人工判断，以及风险减少是否超出增加拒绝所能解释的程度，从而形成有适用边界的监测依据。
